\slajdid{08-s030}
\begin{frame}[label=08-s030]
\gusttekst
Za određivanje realnih maksimalnih vrednosti ovih struja, čestu ulogu ima i koju disipaciju kolo može da “trpi”.\par
Na primer za napon logičke nule na izlazu disipacija na kolu koja potiče od izlazne struje je $P=V_{OL}I_{OL}$, a za napon logičke jedinice na izlazu kola disipacija na kolu koja potiče od izlazne struje je $P=(V_{CC}-V_{OH})(-I_{OH})$.\par
Nemojte da vas zbuni ovaj minus. Posledica je neslaganja referentnog i pravog smera struje.\par
I isto tako ne smemo dozvoliti da napon logičke jedinice na izlazu kola padne do $V_{IH}$ odnosno da napon logičke nule ode $V_{IL}$. Sada se povezuju pojmovi
\[V_{OH\min}\leftrightarrow I_{OH\max}\qquad V_{OL\max}\leftrightarrow I_{OL\max}\]
sa koje god strane gledali. Na primer ako smo definisali koliko želimo da bude $V_{OH\min}$ odredićemo koliko je $I_{OH\max}$ da bi naš kriterijum bio zadovoljen. Ili ako smo na primer zbog disipacije definisali $I_{OH\max}$ videćemo i izračunati koliko je $V_{OH\min}$.\par
Pri ovoj analizi na ulaz logičkog kola postavljamo napone tako da dobijemo najgori slučaj.\par
Odnosno \alert{minimalno $I_{OH\max}$}.
\end{frame}
